\section{Numerical comparisons}
\label{comp:section}

The compact family block is useful if a relaxation has a gap that can be
removed using three-variable moments. This section asks two separate
questions: whether the block closes such gaps, and whether the tested
instance classes exhibit them. The results support the first claim on
constructed instances. They give little evidence for an additional gain
on the benchmark and random instances. All experiments concern root
relaxations; no branch-and-bound implementation was tested.

\subsection{Models and comparison protocol}
\label{comp:protocol}

We minimize $x^THx+g^Tx$ over $C_n$, with lifted objective
$f(m,Y)=\inner{H}{Y}+g^Tm$. For $\bar u=(1,u^T)^T$, we use
$\mathcal M_3=\conv\{\bar u\bar u^T:u\in C_3\}$ for the normalized
$4\times4$ matrix form of $\mathcal Q_3$.
The baseline, denoted $B$, contains
the Shor condition, McCormick inequalities, diagonal caps $Y_{ii}\leq m_i$,
and the four triangle inequalities on each considered triple. For dense
instances the Shor matrix has order $n+1$ and all pairs are recorded.
For sparse instances it is imposed on maximal cliques of a chordal
pattern containing the objective support. The other constraints then use
pattern pairs and triples within those cliques. The comparisons introduce
no first or second moment outside this pattern.

The strengthening methods are as follows.
\begin{description}[leftmargin=2.7em,labelwidth=2.1em]
\item[$K$] The three-positive-variable disjoint-support system of
\citet[eq.~(17)]{Khajavirad2026SparseBoxQP}: eight level-three RLT
inequalities, twelve order-two blocks, and six order-three PSD blocks
per selected triple. Its order-four block is already in the Shor
condition. Higher monomials are shared when the same monomial occurs in
different triples. Larger positive-variable sets and nonempty
negative-variable sets were not tested.
\item[$A$] The eight trilinear inequalities and the switched and permuted
second-order-cone constraints of
\citet[eqs.~(14)--(16)]{AnstreicherPuges2025ETRI}.
\item[$F$] The family block of Section~\ref{family:section}, added only
for a violated triple and orientation. The $24$ orientations choose
the distinguished coordinate and independently complement the three
coordinates.
\item[$X$] The exact three-variable moment lift
\eqref{comp:exact-lift}, using five tetrahedra. A selected triple adds
five $4\times4$ DNN blocks and ten linking equations.
\end{description}
Combined labels mean that the corresponding systems are imposed together.
In particular, $KA$ adds both existing systems, and $KAF$ adds the
family blocks to $KA$. The method $XF$ uses the exact lift but selects
triples by family separation. Individual family cuts and individual
cube-hull cuts were also compared with the blocks.

Family separation minimizes $v^T(B-bb^T)v$ over $v\geq0$,
$\mathbf1^Tv=1$, by the fifteen-support calculation in
Proposition~\ref{comp:family-separation}. For $K$, $A$, $KA$, and $X$,
selection instead uses the full hull-depth problem
\eqref{comp:depth}. Hull depth is the smallest evaluation at the triple
moments of a cube-nonnegative quadratic with $\int_{C_3}q(u)\,du=1$;
it is negative exactly outside $\mathcal M_3$ in exact arithmetic.
This is a stronger test than separation of either
existing inequality system, and its cost is reported as separation time.
Triples with a coordinate within $10^{-6}$ of a box bound are skipped
by this hull-based rule. At an exact bound, Shor and McCormick reduce the
triple to a two-variable problem with an exact baseline
\citep[Theorem~6]{AnstreicherBurer2010}; the numerical
skip is a heuristic extension of that fact.

The usual violation threshold is $10^{-6}$. Triangles are separated
again after each solve. The default limits are $200$ new blocks or
triples per round, $25$ rounds, and $3600$ seconds per method; the
$3000$-triangle chains use $1000$ additions and $15$ rounds, and the
benchmark comparisons use $10$ rounds. A run can also stop when the
relative bound change is below $10^{-7}$ over three rounds. These
limits define selective comparisons; they are not full-system
optimality certificates.

The main solver is Clarabel~0.11.1
\citep{GoulartChen2024,ClarabelSoftware2025} with gap and feasibility tolerances
$10^{-8}$ and one thread. The three-variable comparisons use
$10^{-10}$; the hull-depth subproblems use $10^{-9}$.
SCS~3.3.1 \citep{ODonoghueEtAl2016,SCSSoftware2026}
provides selected accuracy comparisons. Reported relaxation
values use the dual-residual correction in
Proposition~\ref{comp:dual-bound}. We denote these floating-point
lower-bound estimates by $L_M$ for method $M$. They are not interval
certificates. Feasible objective values $U$ come from published benchmark
optima, small-instance face enumeration, or the best archived Gurobi and
local-search incumbents. When $U$ is only a feasible value, the fraction
\begin{equation}\label{comp:closure}
 \frac{L_M-L_B}{U-L_B}
\end{equation}
is a conservative estimate of gap closure, conditional on the accuracy
of the reported lower bounds. It is not a claim that $U$ is optimal.

\subsection{Three-variable gaps}
\label{comp:three-variable-results}

To isolate the effect of a triple block, $109$ objectives were constructed
from separating quadratics of points of $B$ outside
$\mathcal M_3$. The separating quadratics satisfy
$\int_{C_3}q(u)\,du=1$ and have positive coefficients of all three
square terms. Their baseline
gaps have median $0.072$ and range $0.0069$--$0.123$.
The exact lift gives the reference value. All systems are imposed in
full for this comparison, including all $24$ family orientations.

\begin{table}[htbp]
\centering
\caption{Fractions of the three-variable gap closed on $109$ constructed
objectives, computed as $(L_M-L_B)/(L_X-L_B)$.}
\label{comp:hard-three-table}
\begin{tabular}{lrrr}
\toprule
Method & Mean & Median & Minimum\\
\midrule
$K$   & $0.9084$ & $0.9177$ & $0.7668$\\
$A$   & $0.5766$ & $0.5811$ & $0.2619$\\
$KA$  & $0.9084$ & $0.9177$ & $0.7668$\\
$F$   & $1.0000$ & $1.0000$ & $0.999983$\\
$KAF$ & $1.0000$ & $1.0000$ & $0.999912$\\
\bottomrule
\end{tabular}
\end{table}

Every $F$ primal value is at most $5.5\times10^{-9}$ below the
corrected exact-lift reference. The small losses in its corrected-bound
closures in Table~\ref{comp:hard-three-table} also reflect dual-residual
correction; they are not accuracy guarantees at the solver's requested
tolerance. The existing systems leave a visible fraction of the gap.
In separate selective
runs, exactly one family orientation was violated at each baseline
point, and its block closed the gap in one added-block solve. Individual
family cuts had median closure $0.38$ after one round and $0.9933$
after $29$ rounds; none met the stopping test within the $30$-round
limit. The comparison illustrates the difference between imposing the
whole parameterized family and adding its current most violated member.

Additional searches did not find a gap between the full family relaxation
and the exact lift. They included $2180$ perturbations of the constructed
objectives and $256\,556$ sampled baseline points satisfying every
family orientation. The smallest computed hull depth in the latter set
was $-6.4\times10^{-8}$. This evidence is consistent with the
three-variable completeness question, but does not prove it. The samples
are finite, nonuniform, and assessed in floating point.

\subsection{Constructed sparse instances}
\label{comp:constructed}

The same hard objectives were assembled into three sparse families.
A chain has $m$ disjoint triangles, with consecutive triangles joined by
one bridge edge; the bridge scale is denoted $\eta$.
A cactus adds each new triangle by sharing one vertex with an earlier
triangle. The third family puts hard objectives on the triangles of
a $2$-tree or $3$-tree, so triangles share edges. Random weights,
permutations, and coefficient perturbations prevent exact repetition
of a single subproblem. These constructions test the family when local
gaps are deliberately present; they do not model an established
application class.

\begin{table}[htbp]
\centering
\small
\caption{Representative constructed instances, all with seed $1$. Closures use
\eqref{comp:closure}. Solve times are cumulative seconds
after the common baseline solve. $F/X$ counts family blocks and lifted
triples, respectively. In the last two columns, each slash separates
the $F$ value from the $X$ value; it is not a ratio.}
\label{comp:constructed-table}
\begin{tabular}{lrrrrrr}
\toprule
Instance & $n$ & $KA$ & $F$ & $X$ &
Solve $F/X$ & Blocks $F/X$\\
\midrule
Chain, $m=100$, $\eta=0$ & $300$ & $.9956$ & $1.0000$ & $1.0000$
 & $.47/2.17$ & $78/78$\\
Chain, $m=300$, $\eta=.3$ & $900$ & $.8577$ & $.8608$ & $.8597$
 & $2.25/24.62$ & $110/223$\\
Chain, $m=1000$, $\eta=.3$ & $3000$ & $.8985$ & $.9021$ & $.8950$
 & $8.13/110.62$ & $346/997$\\
Chain, $m=3000$, $\eta=.3$ & $9000$ & $.8847$ & $.8889$ & $.8807$
 & $26.20/201.36$ & $1051/2999$\\
Cactus, $m=30$ & $61$ & $.8345$ & $.8718$ & $.8718$
 & $.04/.14$ & $27/27$\\
Cactus, $m=300$ & $601$ & $.7440$ & $.7501$ & $.7489$
 & $2.13/11.71$ & $235/284$\\
Cactus, $m=1000$ & $2001$ & $.7135$ & $.7208$ & $.7181$
 & $19.37/97.19$ & $805/997$\\
Hard $2$-tree, $n=60$ & $60$ & $.0046$ & $.0051$ & $.0041$
 & $.06/.35$ & $4/17$\\
\bottomrule
\end{tabular}
\end{table}

On the $16$ chains and eight cacti, the $F$ and $X$ primal values differ
by at most $5.6\times10^{-6}$ relative. The larger differences between
their corrected closures in Table~\ref{comp:constructed-table} reflect
the accuracy limitations discussed in Section~\ref{comp:numerics};
selected higher-accuracy comparisons support a numerical tie.
The family added
$0.31$--$3.73$ percentage points of closure beyond $KA$ when $KA$
had not already matched the strengthened values. Combining $K$ and $A$ with the
family gave no detectable further gain. With $\eta=0$, the chain
separates into independent three-variable problems, so the triple-level
bound is the true optimum. Bridges and vertex sharing can leave a
global gap after all tested triple strengthenings.

The block sizes explain one possible advantage: a family orientation
adds one $5\times5$ PSD block and six nonnegative auxiliaries, whereas
the exact lift adds five $4\times4$ blocks, thirty off-diagonal
nonnegativity constraints, and ten equations. If all $24$ orientations
are needed, the family can be larger. In the selective runs, however,
$7598$ violated orientation--triple pairs occurred on $7459$ violated
triples, or $1.02$ orientations per triple on average.

Selection also matters. At $m=3000$, $X$ lifts $2999$ triples,
whereas $F$ adds $1051$ blocks. The method $XF$ uses the same family
selection rule as $F$, although its different iterates can select
different triples. On chains, its solve times range from $0.23$ to
$2.05$ times those of $F$, in separate runs. Thus the chain comparison
does not isolate an advantage of the block type. On cacti, $XF$ and
$F$ were run in the same invocation: $XF$ took $2.1$--$5.7$ times
the solve time of $F$. This supports a block-size advantage there,
subject to the timing limitations below. Full records, including
separation time and model size, are provided in the companion archive.

Edge overlap gave a different outcome. Across eight hard $k$-tree
instances, the methods closed at most $0.514\%$ of $U-L_B$, and
the hull test selected only zero to seven triples at the baseline
points. Summing individually hard three-variable objectives therefore
does not ensure that the combined problem retains substantial
three-variable gaps. This is an empirical observation for these
instances, not a theorem about overlapping triangles.

\subsection{Benchmark and random instances}
\label{comp:natural}

The benchmark uses the BoxQP collection \citep{BurerBoxQPInstances2019}.
The collection uses the maximization objective $\tfrac12x^TQx+c^Tx$;
we set $H=-Q/2$ and $g=-c$ and negate its reference optima to obtain
the minimization form used here.
Table~\ref{comp:natural-table} summarizes the tested classes. The sparse
positive-diagonal instances are random $2$-trees and $3$-trees with
positive integer square coefficients and signed integer edge and
linear coefficients. Their clique triples all have three positive
square coefficients.

\begin{table}[htbp]
\centering
\small
\caption{Benchmark and random-instance findings. ``Tight'' refers to
the stated numerical tolerance, not an exact certificate.}
\label{comp:natural-table}
\begin{tabular}{p{.25\linewidth}rp{.53\linewidth}}
\toprule
Class & Count & Finding\\
\midrule
BoxQP benchmark & $99$ & $B$ has relative gap below $10^{-6}$ on
$82$ instances. No family orientation was selected in the $17$
original gap audits.\\
Small dense, BoxQP-style generator & $18\,000$ & $n=5,\ldots,10$;
all baseline gaps below $2.3\times10^{-7}$ relative using corrected
dual estimates.\\
Small dense, AP variant & $12\,000$ & One gap of
$3.6\times10^{-6}$ relative, already closed by $KA$; no additional
family gain.\\
Sparse, positive diagonals & $16$ & $n=100$--$3000$; no family
orientation selected. Computed minimum triple depths at least
$-3\times10^{-7}$.\\
\bottomrule
\end{tabular}
\end{table}

The original BoxQP audits illustrate why a numerical hull violation
must be checked against residuals of the baseline itself.
They evaluated every triple in each of the $17$ instances with a
recorded gap. Four computed depths were below $-10^{-6}$, but these
were explained by triangle inequalities left violated by the original
early stopping rule. A triangle quadratic has uniform mean $1/4$;
a violation $r$ therefore forces normalized hull depth at most $-4r$.
Likewise, a diagonal-cap violation $Y_{ii}-m_i=r>0$ forces depth
at most $-6r$, since $x_i-x_i^2$ has uniform mean $1/6$.
These are errors in satisfying $B$, not evidence of a gap between a
feasible baseline point and the triple hull.

One complete stricter audit is available, for \texttt{spar090-075-1}.
It separates triangles at $10^{-8}$ and evaluates all $117\,480$
triples. The computed minimum depth is $-8.10\times10^{-9}$,
with reported primal infeasibility $2.02\times10^{-9}$. This computed
depth is not a one-sided certificate: a remaining triangle residual
implies that the exact minimum is below $-1.4\times10^{-8}$,
and stored depths exceed exact triangle or cap upper bounds by up
to $9.67\times10^{-8}$.
The gain expression in \eqref{comp:gain}, evaluated at the computed
depths, is $0.045\%$ of the recorded gap, or $0.54\%$ after adding
the primal-to-dual margin. If every computed depth overestimates its
exact value by at most $10^{-7}$, these figures become at most
$0.60\%$ and $1.10\%$, respectively. That assumption and exact
baseline feasibility are not certified. Only this one strict
all-triple audit was completed; its estimate does not apply uniformly
to the other $16$ instances. The remaining evidence is consistent
with small triple-local gains, but does not prove a uniform bound.

Proposition~\ref{comp:gain-bound} does not cover $K$ when higher
moments are shared across triples. A separate test on
\texttt{spar050-050-1} added $K$ to all $19\,600$ triples.
Its primal value changed by $4\times10^{-7}$ against a gap of $1.72$.
This is a negative finding on one instance, not a general exclusion
of gains from higher-order consistency.

The small dense BoxQP-style generator reimplements the published
distribution with different random numbers. An additional generator
was based on the description in
\citet{AnstreicherPuges2025ETRI}, but includes the diagonal in the
density mask. It does not reproduce their twelve reported gap
instances. Its one gap, at $n=9$, is approximately $0.00104$ in
absolute value. A targeted audit finds two violated triples; $K$,
$A$, $KA$, and two family blocks all close the gap to their observed
accuracy, while $KAF$ adds no family block after $KA$.
Consequently, the study supplies a natural-generator triple-level gap,
but no natural instance on which the family adds a detectable gain
beyond the existing systems. These generator-specific findings do not
justify a statement about all random box quadratic programs.

\subsection{Accuracy, timing, and interpretation}
\label{comp:numerics}

The companion archive contains compact tabulations, constructed-instance
data, selected raw records, implementation snapshots, and source checksums.
It supports independent inspection of the reported tables. Reproduction
of all benchmark audits and sampling campaigns also requires the larger
original input and output archives and their software dependencies.

The largest models often returned \texttt{AlmostSolved}. The exact
lift had the largest numerical degradation: each triple adds several
blocks and equations linking small lifted matrices to the original
moments. For example, on the $1000$-triangle chain in
Table~\ref{comp:constructed-table}, Clarabel reports corrected values
$-2502.697859$ for $F$ and $-2502.791210$ for $X$.
At tolerance $10^{-6}$, SCS instead gives $-2502.697709$ and
$-2502.697835$, respectively. Their relative difference is
$5.05\times10^{-8}$, supporting a numerical tie. The lower
Clarabel value for $X$ is therefore an accuracy limitation, not
evidence that the family is mathematically stronger than the exact
lift. Primal-to-dual margins are retained in the records and must be
considered before interpreting small differences.

Times are wall clock on a shared $36$-core machine, whose load varied
substantially. Methods ran sequentially, so even within-run ratios
are affected by changing load. A separate quieter run of the
$300$-triangle chain reproduced the bounds but changed $X/F$
from $10.9$ to $7.5$ and $KA/F$ from $3.2$ to $1.4$.
The tabulated solve times exclude the common baseline stage and
separation time; they do not measure complete solver time.
The model also stores the objective matrix densely, limiting the
largest tested problem to $9000$ variables. These observations
support a useful compact strengthening on deliberately constructed
triple gaps. They establish neither a general runtime advantage nor
an improvement in a global optimization solver.
